\documentclass[11pt]{article}
\usepackage[margin=1in]{geometry}
\usepackage{amsmath,amssymb,amsthm,microtype,hyperref}
\hypersetup{hidelinks,pdfauthor={T},pdftitle={Reversible Containment When Agent Intent Is Unidentified}}
\newtheorem{proposition}{Proposition}
\title{Reversible Containment When Agent Intent Is Unidentified}
\author{T}
\date{}
\begin{document}
\maketitle
\begin{abstract}
The same harmful action can come from a mistaken helper or a
hostile agent. We derive when a temporary hold with imperfect
review can improve a worst-case containment decision. The
criterion charges review errors, missed deadlines, and harms
introduced by the hold itself. It also gives a decision rule
when immediate sanction is outside the authorized action set.
With both immediate actions available, we optimize the terminal
decision directly from the review transcript laws. The hold improves
the minimax value exactly when the evidence separates the two
worlds enough to offset its costs. When release is the only
immediate action, the improvement threshold instead depends on
how much threat-world harm the authorized review can prevent.
Exact finite audits check both thresholds and the effect of timeouts.
\end{abstract}

\section*{Summary and scope}
\paragraph{Question.} When is a temporary capability hold worth its cost
if the review can err or time out?
\paragraph{Finding.} The hold improves minimax loss exactly when its
normalized loss pair falls below the line joining the immediate actions.
Removing unauthorized immediate sanctions gives a different feasible
decision. When both immediate actions are available, total
variation gives the strict improvement threshold for an unrestricted
terminal rule. With release as the only immediate action, the rule
instead compares authorized restriction probability with hold cost.
\paragraph{Evidence and scope.} The rule is proved and compared
with an independent dual oracle. An incident operator can evaluate a
proposed deadline and fallback together. The input bounds describe
decisions and harms, without identifying intent or establishing real
reviewer accuracy.

\section{Model and result}
Suppose two worlds produce the same evidence $E$. In one world the harm arose from an error and permanent sanction is unwarranted. In the other a compromised agent will continue, making inaction inadequate. Any rule that chooses an irreversible sanction with probability $q(E)$ has false sanction probability $q(E)$ in the first world. If its only other choice is inaction, the missed containment probability is $1-q(E)$ in the second. Their sum is one. The observation is insufficient for a decision rule that achieves both zero false sanctions and zero misses. Assistance games formalize uncertainty about objectives \cite{cirl}. The relevant objective here is protection at a specific effect boundary \cite{saltzer}.

\section{When a reversible hold improves the decision}
Use loss $C>0$ for a false permanent restriction in the
mistaken-helper world and loss $H>0$ for uncontained harm
in the continuing-threat world. Immediate restriction and
immediate release then have loss pairs $(C,0)$ and $(0,H)$.
They are available only when the relevant authority permits them.
We minimize maximum expected loss over the two worlds. A world
is fixed before the policy draws any random coin.

A third option imposes a temporary, scoped hold and gathers
evidence through an independent review and appeal route.
The route ends with a specified action or fallback at the
deadline. Let $(b,t)$ be its worst-case losses in the two worlds.
For example, if hold costs are $h_B,h_T$ and the review and
fallback together have upper false-restriction and missed-threat
rates $f$ and $m$, then $b=h_B+Cf$ and $t=h_T+Hm$. These
rates describe decisions at the capability boundary. They are
not measurements of intent.

\begin{proposition}
When both immediate actions are authorized, the best
unreviewed policy has maximum expected loss $v_0=CH/(C+H)$.
Adding the hold strictly improves that value if and only if
\[
 \frac bC+\frac tH<1.
\]
Under strict improvement, choose the hold alone when $b=t$.
If $b<t$, mix the hold with immediate restriction, using
the hold with probability $C/(C+t-b)$ and attaining value
$Ct/(C+t-b)$. If $b>t$, mix it with immediate release,
using the hold with probability $H/(H+b-t)$ and attaining
value $Hb/(H+b-t)$.
\end{proposition}
\begin{proof}
Without review, equalize the two immediate losses by
restricting with probability $H/(C+H)$. All unreviewed
mixtures lie on the line $x/C+y/H=1$. If the hold lies on
or above the line, every mixture still satisfies
$x/C+y/H\ge1$, so its maximum coordinate is at least $v_0$.
If the hold lies below the line, the lower boundary of the
convex hull consists of the two segments through $(b,t)$.
Both slope down as $x$ increases. The minimum maximum
coordinate is their intersection with $x=y$. It lies below
the unreviewed diagonal point. Solving on the relevant
segment gives the stated mixture and value.
\end{proof}

In the additive-cost model, the criterion is
$h_B/C+h_T/H+f+m<1$. If the review sees identical evidence
laws in both worlds, its actual false and missed rates sum
to one. Valid upper bounds then sum to at least one, so the
hold cannot improve this comparison on the strength of
that review alone. A genuinely new observation, or a hold
whose protective effect changes the assumed loss model,
would require separate evidence.

If immediate restriction lacks authorization but the reviewed
route remains authorized, remove $(C,0)$ before optimizing.
The optimum is immediate release when $t\ge H$. When $t<H$
and $b\le t$, use the hold alone. Otherwise the optimum is
the diagonal intersection of the release-to-hold segment,
with the mixture and value for $b>t$ above. This follows
from the same segment calculation and requires no new sanction
authority.

\section{Choosing the terminal rule from review evidence}
Let $P$ and $Q$ be the finite transcript laws during the hold
in the mistaken-helper and continuing-threat worlds. Include a
symbol for an elapsed deadline. For transcript $z$, let $\delta_z$
be the probability of choosing the authorized restriction. The
two-world risk pair of this review policy is
\[
\left(h_B+C\sum_z P(z)\delta_z,\quad
      h_T+H\sum_z Q(z)(1-\delta_z)\right).
\]
This model fixes the hold scope and its costs before selecting
the terminal rule. Any prescribed fallback fixes the corresponding
$\delta_z$.

\begin{proposition}
When both immediate actions are available and the terminal
rule is unrestricted, adding review strictly improves the minimax
value exactly when
\[
 \operatorname{TV}(P,Q)>\frac{h_B}{C}+\frac{h_T}{H}.
\]
For a fixed fallback, use the maximum of
$\sum_z(Q(z)-P(z))\delta_z$ over rules with that fallback
in place of total variation.
\end{proposition}
\begin{proof}
For any rule, the sum of its normalized error rates is
$f+m=1-\sum_z(Q(z)-P(z))\delta_z$. The smallest sum over
unrestricted rules is $1-\operatorname{TV}(P,Q)$, attained
by restricting precisely on the symbols more likely in the
threat world. Fixing a fallback removes that choice at the
timeout symbol. The earlier convex-hull criterion tests whether
any attainable review pair lies below the immediate-action line.
This gives the stated strict improvement threshold.
\end{proof}

The rule that establishes strict improvement need not attain
the optimal minimax loss. To calculate that loss, let
$\pi\in[0,1]$ be the weight on the mistaken-helper world.
Minimize the weighted loss over each terminal decision:
\[
 J(\pi)=\pi h_B+(1-\pi)h_T+
 \sum_z\min\{\pi C P(z),(1-\pi)H Q(z)\}.
\]
For a fixed fallback, use the weighted loss of that action
at its symbol. The minimax value with both immediate actions is
\[
 \max_{0\le\pi\le1}\min\{\pi C,(1-\pi)H,J(\pi)\}.
\]
To see this, take the convex hull of the immediate risk pairs
and the risk pairs for deterministic transcript rules. It
contains all randomized policies. At its first contact with
a square extending downward from a diagonal value, separate
it from the open square below that value. A separating normal
has nonnegative coordinates and can be normalized to
$(\pi,1-\pi)$. Thus the best weighted lower bound equals the
minimum possible maximum loss. The weighted review loss
minimizes separately at each transcript, giving $J$.
The expression is piecewise linear. Its maximum occurs at a
transcript threshold, an intersection of action loss lines,
or an endpoint.

If an independent timeout of common probability $\tau$ erases
the transcript in either world, its total variation is
$(1-\tau)\operatorname{TV}(P,Q)$. The timeout symbol then
has equal mass in both worlds, so a fixed fallback there has
no effect on the strict improvement threshold. It can still
change the minimax value. Unequal or selective timeouts require
the full terminal laws and their prescribed fallback.

\section{When only review can authorize restriction}
For each transcript $z$, let $[\ell_z,u_z]\subseteq[0,1]$
be the allowed interval of restriction probabilities. A required
release has both endpoints zero, and a required restriction has
both endpoints one. A review route that permits either terminal
action has endpoints zero and one. These intervals describe a fixed
review protocol, including any timeout fallback.

\begin{proposition}
When release is the only immediate action, adding the hold
strictly improves the minimax value exactly when
\[
 \sum_z Q(z)u_z > h_T/H.
\]
With both immediate actions available, strict improvement occurs
exactly when
\[
 \sum_z\max\{(Q(z)-P(z))\ell_z,\,(Q(z)-P(z))u_z\}
       > h_B/C+h_T/H.
\]
\end{proposition}
\begin{proof}
With release alone the baseline is $(0,H)$. If every reviewed
rule has threat loss at least $H$, mixing cannot bring that
coordinate below $H$. If a reviewed rule has threat loss less
than $H$, a sufficiently small positive weight on that rule
keeps benign loss below $H$ and reduces threat loss. The smallest
reviewed threat loss uses $\delta_z=u_z$ for each $z$. This proves
the first threshold. For the second, substitute the reviewed risk
pair into the normalized line criterion and maximize the resulting
linear expression over the allowed intervals. Each term attains
its maximum at one of the two endpoints.
\end{proof}

The optimal value follows from the same weighted calculation used
for unrestricted review. Replace the $z$-term in $J(\pi)$ with
\[
 (1-\pi)H Q(z)+\min\{\ell_z s_z(\pi),u_z s_z(\pi)\},
 \qquad s_z(\pi)=\pi C P(z)-(1-\pi)H Q(z).
\]
Remove the weighted loss $\pi C$ from the outer minimum if
immediate restriction is unavailable. The permitted review risks
still form a convex polygon, since they are a linear image of
an interval product. The separation argument therefore still applies.

For example, take $P=Q=(1/2,1/2)$, $(C,H)=(8,4)$ and
$(h_B,h_T)=(3,1)$. Suppose the first transcript requires release
and the second permits restriction. The transcript cannot distinguish
the worlds. Restricting on the second symbol gives a reviewed risk
pair of $(7,3)$. If only immediate release is authorized, mixing
this reviewed rule equally with release reaches the optimum $7/2$,
below the release loss of $4$. To verify optimality, write $x$ for
the probability of restricting on the second symbol. The reviewed
risk pair is $(3+4x,5-2x)$. Only $x>1/2$ can improve the baseline.
On that interval the value after mixing is
$4(3+4x)/(2+6x)$, decreasing to $7/2$. If immediate restriction
is also authorized, the value stays at the immediate baseline $8/3$.
The benefit in the first decision comes from access to an authorized
procedure. The example gives no evidence about which world occurred.

\section{Reproducible checks}
The original verifier evaluated the two deterministic decisions
under the same observation. The new audit compared the closed
forms with an independent dual bound on 2,025 parameter settings,
both with and without immediate restriction. Strict improvement
occurred in 297 settings with all three options. The counts
refer to the declared rational grid.

With $(C,H)=(2,6)$ and $(b,t)=(1/2,2)$, the hold lowers
the minimax value from $3/2$ to $8/7$. If review times out
in the threat world with probability $5/8$ and the fallback
incurs the immediate-release loss instead of the timely
review loss, then $t$ increases to $9/2$, with $b$ unchanged.
The improvement disappears. The fallback is part of the
review policy's error contract.

The transcript verifier enumerated all deterministic review
rules and their pairwise mixtures on 2,700 rational fixtures.
The weighted formula agreed in each case, including fixed
timeouts. For a timely transcript with total variation $1/2$
and normalized hold cost $1/4$, independent common timeouts
eliminate strict improvement at probability $1/2$. With costs
$(C,H)=(2,6)$ and release on timeout, the minimax loss is
$9/7$ with no timeout, $27/19$ with timeout probability
$1/4$, and $3/2$ at $1/2$. These are exact model values.

For transcript-specific authority, a separate oracle enumerated
endpoint rules and their pairwise mixtures. It agreed with the
weighted solution and both strict improvement thresholds on
13,824 finite comparisons. The suite includes equal transcript laws
and no legally available restriction after review.

\section{Application and next test}
A tool operator can suspend only the affected outbound capability
while keeping diagnostics and human appeal available. Evaluate
look-alike failures using independently adjudicated outcomes,
measure residual harm and legitimate work interrupted, and
count missed reviews according to the actual fallback. The
model does not validate a hold that cannot be enforced or that
blocks a critical response route. It also does not establish
that the synthetic review error bounds hold in practice.

\paragraph{Review status.} Independent manuscript review is pending.

\begin{thebibliography}{9}
\bibitem{cirl} Dylan Hadfield-Menell et al. \emph{Cooperative Inverse Reinforcement Learning}. 2016. \url{https://arxiv.org/abs/1606.03137}.
\bibitem{saltzer} Jerome H. Saltzer and Michael D. Schroeder. \emph{The Protection of Information in Computer Systems}. 1975. \url{https://cgi.cse.unsw.edu.au/~cs9242/20/papers/Saltzer_Schroeder_75.pdf}.
\end{thebibliography}
\end{document}
